%\part{First-order Numerical Methods}
\chapter{Asymptotic Compatibility (AC) of the Numerical Methods}
In this chapter, we demonstrate the convergence analysis of the numerical solution $\mathcal{U}^{\epsilon, \Delta x}$ to the unique entropy solution of the local model \eqref{eq:local model} as both $\Delta x \rightarrow 0$ and $\epsilon \rightarrow 0$ simultaneously.
\section{Convergence Properties of Numerical Schemes for the Local Model}
Before studying the convergence of numerical methods for the nonlocal model to local model, it is significant to study the convergence properties of numerical methods for the local model. Here is a thoughtful question to ask: if we have convergence criteria (or properties) for the numerical methods of the local model, do the numerical methods of the nonlocal model satisfy those criteria?
\subsection{Conservative Schemes}
The first property we require of the numerical schemes for the local model is conservation: no quantity should be created or destroyed over time. The entropy solution of the local model is conservative, in the sense that the integral of the solution is preserved over time. Therefore, we require the corresponding discrete quantity to be preserved as well.
\begin{definition}[Conservative scheme]
    The three-point update function $H$ given in \eqref{eq:three_point_H} is conservative if satisfies
    \[ \sum_{j}^{} u_{j}^{n+1} = \sum_{j}^{} u_{j}^{n} \quad \text{ for all } n .\]
\end{definition}
\subsection{Consistent Schemes}
The numerical schemes must be consistent with the local model \eqref{eq:local model}, i.e, they must approximate the desired PDE.
\begin{definition}[Consistency]
    The three-point update function $H$ with corresponding flux $F$ is consistent if
    \[ F(w,\cdots, w) = f(w) \quad \text{for all } w \in \mathbb{R}.\]
\end{definition}
\subsection{Monotone Schemes}
\begin{definition}[Monotonicity]
    The three-point update function $H$ is non-decreasing in each of its arguments.
\end{definition}
\subsection{$L^{\infty}$ Bounds}
\subsection{Entropy Inequalities and $L^{p}$ Bounds}
\subsection{TV Bounds}
\subsection{Time Continuity}
\subsection{Helly's Theorem}
\subsection{Ascoli's Theorem}
\subsection{Lax--Wendroff Theorem}
\section{Convergence of Numerical Schemes for the Nonlocal Model to Local model}
For the local model, if we have a conservative, consistent and monotone scheme, then the scheme converges to the correct solution. However, for the nonlocal model, it is difficult to obtain monotone scheme. To see this, we consider the update function $H$ given in \eqref{eq:generic_H_N}. In \textcolor{red}{Citation is needed here}, they have proved the Asymptotically compatibility theorem for general fux $g$. To make the proof easier to follow up, we restrict ourselves to the Godunov flux $g$, and provide the proof in this chapter. We start making the necessary assumptions on the numerical quadrature weights $w_k$, initial data $\rho_0$, Godunov numerical flux $g$ and CFL $\lambda = \frac{\Delta t}{\Delta x}$
\begin{assumption}
    \leavevmode
    The nonlocal kernel is given by
    \[
        w_\epsilon(y) = \frac{1}{\epsilon} w\!\left(\frac{y}{\epsilon}\right),
        \qquad y \in [0,\epsilon],
    \]
    where $w$ is a $C^1$ nonnegative, strictly decreasing probability density on $[0,1]$, satisfying
    \[
        \int_0^1 w(y)\,dy = 1.
    \]
\end{assumption}
\begin{assumption}
    The initial data $\rho_0$ belongs to $L^\infty(\mathbb{R})$, has bounded total variation, and satisfies the following:
    \begin{enumerate}
        \item
              \[
                  0 \le \rho_0(x) \le 1
                  \qquad \text{for all } x \in \mathbb{R}.
              \]
        \item One-sided Lipschitz condition
              \[
                  \rho_0(x) \ge \rho_{\min} > 0
                  \qquad \text{for all } x \in \mathbb{R},  \qquad  -\frac{\rho_0(y)-\rho_0(x)}{y-x} \le L
                  \qquad \text{for all } x \ne y \in \mathbb{R},
              \]
              for some constants $\rho_{\min}>0$ and $L>0$.
    \end{enumerate}
\end{assumption}

\begin{assumption}
    The numerical quadrature weights $\{w_k\}_{0 \le k \le m-1}$ satisfy
    \[
        \omega_{\epsilon}(k\Delta x) \Delta x \ge w_k \ge \omega_{\epsilon}((k+1)\Delta x) \Delta x
    \]
    and
    \[
        w_k - w_{k+1} \ge c\,m^{-2},
    \]
    for some constant $c>0$ depending only on the kernel $w$.
    Moreover, the weights satisfy the normalization condition
    \[
        \sum_{k=0}^{m-1} w_k = 1.
    \]
\end{assumption}
\begin{assumption}
    \leavevmode
    \begin{enumerate}
        \item The numerical flux function $g$ is quadratic.

        \item When $\rho_L=\rho_R$ and $q_L=q_R$, one has
              \[
                  g(\rho_L,\rho_L,q_L,q_L) = \rho_L(1-q_L).
              \]

        \item Denote by $\gamma_{ij}$, $1 \le i,j \le 4$, the second-order partial derivatives of $g$. Then
              \[
                  \gamma_{11}=\gamma_{12}=\gamma_{22}=0,
                  \qquad
                  \gamma_{33}=\gamma_{34}=\gamma_{44}=0,
              \]
              and
              \[
                  \gamma_{13},\gamma_{23},\gamma_{14},\gamma_{24} \le 0,
                  \qquad
                  \gamma_{13}+\gamma_{23}+\gamma_{14}+\gamma_{24} = -1.
              \]
    \end{enumerate}
\end{assumption}
\begin{assumption}
    Denote by $\theta^{(i)}$, $1 \le i \le 4$, the first-order partial derivatives of $g$
    with respect to the four arguments $\rho_L,\rho_R,q_L,q_R$. For any
    \[
        0 \le \rho_L,\rho_R,q_L,q_R \le 1,
    \]
    we assume the following:
    \begin{enumerate}
        \item
              \begin{align*}
                   & \theta^{(1)}(q_L,q_R) \ge 0,
                  \qquad
                  \theta^{(2)}(q_L,q_R) \le 0,          \\
                   & \theta^{(3)}(\rho_L,\rho_R) \le 0,
                  \qquad
                  \theta^{(4)}(\rho_L,\rho_R) \le 0.
              \end{align*}

        \item
              \begin{align*}
                   & \theta^{(1)}(q_L,q_R)+\theta^{(3)}(\rho_L,\rho_R)
                  +2(\gamma_{13}+\gamma_{23}) \ge 0,                         \\
                   & \theta^{(2)}(q_L,q_R)-2(\gamma_{23}+\gamma_{24}) \le 0.
              \end{align*}

        \item
              \[
                  \theta^{(3)}(\rho_L,\rho_R)+\theta^{(4)}(\rho_L,\rho_R)
                  \le -\min\{\rho_L,\rho_R\}.
              \]
    \end{enumerate}
\end{assumption}

\begin{assumption}
    With the notation $\theta^{(i)}$, $1 \le i \le 4$, from the previous assumption, $\lambda$ satisfies
    \[
        \lambda \sum_{i=1}^4 \|\theta^{(i)}\|_{L^{\infty}} < 1.
    \]
\end{assumption}
Let us now check if the Godunov flux $g$ satisfy the Assumptions 4.2.4 and 4.2.5. We begin by computing the first order and second order partial derivatives of $g$ with respect to the four arguments $\rho_L,\rho_R,q_L,q_R$.
The Godunov numerical flux is given by
\[g(\rho_L,\rho_R,q_L,q_R)= \rho_L (1-q_R),\] and we compute the partial derivatives as follows:

\begin{align*}
    \theta^{(1)}(q_L,q_R) & := \frac{\partial g}{\partial \rho_L}, & \theta^{(2)}(q_L,q_R) & := \frac{\partial g}{\partial \rho_R} \\
                          & = 1-q_R                                &                       & = 0
\end{align*}
\begin{align*}
    \theta^{(3)}(\rho_L,\rho_R) & := \frac{\partial g}{\partial q_L}, & \theta^{(4)}(\rho_L,\rho_R) & := \frac{\partial g}{\partial q_R} \\
                                & = 0                                 &                             & = -\rho_L.
\end{align*}
Since $\theta^{(2)}$ and $\theta^{(3)}$ are zero, and  $\theta^{(1)}$ depends only on $q_R$ while $\theta^{(4)}$ on $\rho_L$, the only nonzero second order partial derivatives are the following
\begin{align*}
    \gamma_{14} & = \frac{\partial^2 g}{ \partial q_R \partial \rho_L}, & \gamma_{41} & = \frac{\partial^2 g}{\partial \rho_L \partial q_R} \\
                & = -1                                                  &             & = -1.
\end{align*}
Checking Assumptions 4.2.4 and 4.2.5 is straightforward now, and the Godunov numerical flux satisfies both Assumptions. let us now see what does the CFL $\lambda$ satisfy explicitly, provided that $\theta^{(i)}$ are the computed first order partial derivatives that corresponding to the Godunov numerical flux $g$. Given
\[
    0 \le \rho_L,\rho_R,q_L,q_R \le 1,
\] as Assumption 4.2.4, we obtain
\begin{equation}\label{eq:lambda_God}
    \begin{split}
        \lambda \sum_{i=1}^4 \|\theta^{(i)}\|_{L^{\infty}} < 1                                                                                              \\
        \lambda \Bigl(\|\theta^{(1)}\|_{L^{\infty}}+ \|\theta^{(2)}\|_{L^{\infty}} + \|\theta^{(3)}\|_{L^{\infty}} + \|\theta^{(4)}\|_{L^{\infty}}\Bigr)< 1 \\
        \lambda \Bigl(\|\theta^{(1)}\|_{L^{\infty}}+ 0 + 0 + \|\theta^{(4)}\|_{L^{\infty}}\Bigr)< 1                                                         \\
        \lambda \Bigl(\|\theta^{(1)}\|_{L^{\infty}}+ \|\theta^{(4)}\|_{L^{\infty}}\Bigr)< 1                                                                 \\
        \lambda \Bigl( \sup\limits_{q_R \in [0,1]}| 1-q_R\vert  +  \sup\limits_{\rho_L \in [0,1]}|-\rho_L\vert \Bigr)< 1                                    \\
        2\lambda < 1
    \end{split}
\end{equation}
So, for the Godunov numerical flux $g$, the CFL condition is given by
\[ \lambda < \frac{1}{2}.\]
\subsection{Maximum Principle}
It is easy to establish the maximum principle if we have monotone numerical schemes. let us now check if the update function $\mathcal{H}$ with corresponding Godunov numerical flux $g$ is monotone. To do that, suppose that
\[0 \leq \rho_{\min} \leq  \rho_{j+k}^{n} \leq \rho_{\max} \leq 1 \qquad \text{ for } k = -1, 0,1,\cdots,m,\]
then we have also
\begin{align*}
    q_{j}^{n} & = \sum_{k=0}^{m-1} w_k \rho_{j+k}^{n}              \\
              & \leq \sum_{k=0}^{m-1} w_k \rho_{\max}              \\
              & =  \rho_{\max} \sum_{k=0}^{m-1} w_k = \rho_{\max},
\end{align*}
where we have used $\sum_{k=0}^{m-1} w_k = 1$ by the normalization condition. By similar argument $q_{j}^{n}$ is lower bonded by $\rho_{\min}$. Consequently, we have
\begin{equation}\label{eq:q_bound}
    0 \leq q_{j}^{n} \leq 1 \text{ for all } j\in \mathbb{Z}, n\geq 0
\end{equation}

We check now whether $\mathcal{H}$ is nondecreasing with respect to each of its arguments. The update function $\mathcal{H}$ with respect to the Godunov numerical flux
\[g(\rho_L,\rho_R,q_L,q_R)= \rho_L (1-q_R),\]
is given by
\begin{equation}\label{eq:mono_H}
    \mathcal{H}( \rho_{j-1}^n, \rho_{j}^n,  \rho_{j+1}^n,\cdots,  \rho_{j+m}^n)
    =
    \rho_j^n
    +
    \lambda
    \left[
        \rho_{j-1}^n(1-q_j^n)
        -
        \rho_j^n(1-q_{j+1}^n)
        \right].
\end{equation}
Taking the partial derivatives gives us
\begin{align*}
    \frac{\partial \mathcal{H}}{\partial \rho_{j-1}^n} & =\lambda(1-q_j^n)                                                                           \\
    \frac{\partial \mathcal{H}}{\partial \rho_j^n}     & =1-\lambda\left[(1-q_{j+1}^n)+w_0\rho_{j-1}^n\right]                                        \\
    \frac{\partial \mathcal{H}}{\partial \rho_{j+k}^n} & = \lambda\left( w_{k-1}\rho_j^n -w_k\rho_{j-1}^n\right) \qquad \text{for } k= 1,2,\cdots,m,
\end{align*}
where we make the convention $w_m=0$.

$\mathcal{H}$ is nondecreasing with respect to $\rho_{j-1}^n$, because
\begin{equation}\label{eq:First_argument}
    \begin{split}
        \frac{\partial \mathcal{H}}{\partial \rho_{j-1}^n} & =\lambda(1-q_j^n) \geq 0 \qquad \text{due to } \eqref{eq:q_bound}.
    \end{split}
\end{equation}
Next, we check if $\mathcal{H}$ is nondecreasing with respect to $\rho_{j}^n$. We desire to have that
\[
    \frac{\partial \mathcal{H}}{\partial \rho_j^n}
    =
    1-\lambda\left[(1-q_{j+1}^n)+w_0\rho_{j-1}^n\right] \geq 0.
\]
For this to happen, we must have
\[
    0 \leq \lambda\left[(1-q_{j+1}^n)+w_0\rho_{j-1}^n\right]\leq 1
\]
Clearly, we have $\lambda\left[(1-q_{j+1}^n)+w_0\rho_{j-1}^n\right] \geq 0$, so it remains to show that
\[\lambda\left[(1-q_{j+1}^n)+w_0\rho_{j-1}^n\right] \leq 1\].

We remember that the first order partial derivatives of $g$ are
\begin{align*}
    \theta^{(1)}(q_L,q_R) & = \frac{\partial g}{\partial \rho_L}, & \theta^{(2)}(q_L,q_R) & = \frac{\partial g}{\partial \rho_R} \\
                          & = 1-q_R                               &                       & = 0
\end{align*}
\begin{align*}
    \theta^{(3)}(\rho_L,\rho_R) & = \frac{\partial g}{\partial q_L}, & \theta^{(4)}(\rho_L,\rho_R) & = \frac{\partial g}{\partial q_R} \\
                                & = 0                                &                             & = -\rho_L.
\end{align*}
That means
\begin{align*}
    \lambda\left[(1-q_{j+1}^n)+w_0\rho_{j-1}^n\right] & \leq \lambda\left[(1-q_{j+1}^n)+\rho_{j-1}^n\right]               \qquad \text{ since $0\leq w_0 \leq 1$ and $\rho_{j-1}^n \geq 0$ } \\
                                                      & = \lambda \Bigl(\theta^{(1)}(q_{j}^{n}, q_{j+1}^{n}) - \theta^{(4)}(\rho_{j-1}^{n},\rho_{j}^{n})\Bigr)                               \\
                                                      & \leq \lambda \Bigl(\|\theta^{(1)}\|_{L^{\infty}}+ \|\theta^{(4)}\|_{L^{\infty}}\Bigr)                                                \\
                                                      & = 2\lambda                                                                                                                           \\
                                                      & < 1 \qquad \text{ By Assumption 4.2.4}
\end{align*}
It follows that
\begin{equation}\label{eq:Second_argument}
    \frac{\partial \mathcal{H}}{\partial \rho_j^n}
    =
    1-\lambda\left[(1-q_{j+1}^n)+w_0\rho_{j-1}^n\right] \geq 0.
\end{equation}
Therefore, $\mathcal{H}$ is nondecreasing with respect to $\rho_{j}^n$.

We consider now for $k= 1,2,\cdots, m-1$.
\begin{align*}
    \frac{\partial \mathcal{H}}{\partial \rho_{j+k}^n} & = \lambda\left( w_{k-1}\rho_j^n -w_k\rho_{j-1}^n\right)                      \\
                                                       & =\lambda \Bigl[(w_{k-1}- w_k) \rho_j^n + w_k(\rho_j^n - \rho_{j-1}^n )\Bigr]
\end{align*}
By Assumption 4.2.3
\[(w_{k-1}- w_k) \geq cm^{-2} \geq 0. \] Hence the first term inside the brackets is nonnegative. However, the second term inside the brackets may be negative, because Assumptions 4.2.1--4.2.6 alone do not imply that $\rho_{j-1}^n >\rho_j^n$. Consequently, we may have $\frac{\partial \mathcal{H}}{\partial \rho_{j+k}^n} < 0$.

For $k= m$, the convection $w_m=0$ gives
\begin{align*}
    \frac{\partial \mathcal{H}}{\partial \rho_{j+m}^n} & = \lambda\left( w_{k-1}\rho_j^n -w_k\rho_{j-1}^n\right) \\
                                                       & =\lambda (w_{m-1}\rho_j^n) \geq 0
\end{align*}
Thus, we can not conclude $\mathcal{H}$ is nondecreasing in each of its arguments. Therefore, one can not prove the maximum principle by showing $\mathcal{H}$ is monotone. Instead, we prove the maximum principle by applying the mean value theorem.

\begin{theorem}[Maximum principle]
    Suppose Assumptions  4.2.1--4.2.6 holds. In addition, the following condition is satisfied
    \[ 0 < \epsilon \leq \epsilon_0:= \frac{c\rho_{\min}}{2Lw(0)}\]
    where the constants $\rho_{\min}$ and $L$ are as in Assumption 4.2.2, while the constant $c$ is as in Assumption 4.2.3.  Let $\{\rho_j^n\}_{j \in \mathbb{Z}}^{ n \in \mathbb{N}}$ be the numerical solution computed by \eqref{eq:mono_H}. Then the following uniform bound hold
    \[\rho_{\min} \leq \rho_{j}^{n} \leq  \rho_{\max}  \qquad \text{for all } j \in \mathbb{Z},  n \in \mathbb{N} \]
    where
    \[
        \rho_{\min} := \inf\limits_{x\in \mathbb{R}} \rho_0 > 0 \quad \text{and }  \rho_{\max}:= \sup\limits_{x\in \mathbb{R}} \rho_0 \leq 1
    \]
\end{theorem}

We will prove the above theorem by induction, but first we prove the following Lemma.

\begin{lemma}
    Suppose all conditions in Theorem 4.2.7 are given, and that
    \[0 \leq \rho_{\min} \leq  \rho_{j+k}^{n} \leq \rho_{\max} \leq 1 \qquad \text{ for } k = -1, 0,1,\cdots,m.\]
    Then we have
    \[ \rho_{\min} \leq \mathcal{H}( \rho_{j-1}^n, \rho_{j}^n,  \rho_{j+1}^n,\cdots,  \rho_{j+m}^n) \leq \rho_{\max} \]
\end{lemma}
\begin{proof}
    Since
    \[
        \mathcal{H}(\rho_{\min},\rho_{\min},\rho_{\min},\ldots,\rho_{\min})
        =\rho_{\min},
    \]
    we write
    \[
        \mathcal{H}(\rho_{j-1}^n,\rho_j^n,\rho_{j+1}^n,\ldots,\rho_{j+m}^n)-\rho_{\min}
        =
        \Delta \mathcal{H}_1+\Delta \mathcal{H}_2,
    \]
    where $\Delta \mathcal{H}_1$ and $\Delta \mathcal{H}_2$ are given by
    \[\Delta \mathcal{H}_1 = \mathcal{H}( \rho_{j-1}^n, \rho_{j}^n,  \rho_{j+1}^n,\cdots,  \rho_{j+m}^n) - \mathcal{H}( \rho_{\min}, \rho_{\min},  \rho_{j+1}^n,\cdots,  \rho_{j+m}^n),\]
    \[\Delta \mathcal{H}_2 = \mathcal{H}( \rho_{\min}, \rho_{\min},  \rho_{j+1}^n,\cdots,  \rho_{j+m}^n) - \mathcal{H}(\rho_{\min}, \rho_{\min}, \rho_{\min},\cdots, \rho_{\min})\]

    Focusing on $\Delta \mathcal{H}_1$, by mean value theorem there exists a point $(\tilde{\rho}_{j-1}^n, \tilde{\rho}_{j}^n)$ on the line segment from $(\rho_{\min},\rho_{\min})$ to $(\rho_{j-1}^n,\rho_{j}^n)$ such that
    \begin{align*}
        \Delta \mathcal{H}_1 = \frac{\partial \mathcal{H}}{\partial \rho_{j-1}^n} ( \tilde{\rho}_{j-1}^n, \tilde{\rho}_{j}^n, \rho_{j+1}^n, \cdots, \rho_{j+m}^n)(\rho_{j-1}^n- \rho_{\min}) + \frac{\partial \mathcal{H}}{\partial \rho_{j}^n} ( \tilde{\rho}_{j-1}^n, \tilde{\rho}_{j}^n, \rho_{j+1}^n, \cdots, \rho_{j+m}^n)(\rho_{j}^n- \rho_{\min}),
    \end{align*}
    with
    \[0 \leq \rho_{\min} \leq \tilde{\rho}_{j-1}^n \leq \rho_{\max} \text{ and } 0 \leq \rho_{\min} \leq \tilde{\rho}_{j}^n \leq \rho_{\max}.\]
    We want to show that
    \[
        \frac{\partial \mathcal{H}}{\partial \rho_{j-1}^n} ( \tilde{\rho}_{j-1}^n, \tilde{\rho}_{j}^n, \rho_{j+1}^n, \cdots, \rho_{j+m}^n)(\rho_{j-1}^n- \rho_{\min}) \geq 0
    \]
    and
    \[
        \frac{\partial \mathcal{H}}{\partial \rho_{j}^n} ( \tilde{\rho}_{j-1}^n, \tilde{\rho}_{j}^n, \rho_{j+1}^n, \cdots, \rho_{j+m}^n)(\rho_{j}^n- \rho_{\min}) \geq 0,
    \]
    so that $\Delta \mathcal{H}_1 \geq 0.$
    Since
    \[
        (\rho_{j-1}^n- \rho_{\min}) \geq 0,
        \quad (\rho_{j}^n- \rho_{\min}) \geq 0,
    \]
    it suffices to show that
    \[
        \frac{\partial \mathcal{H}}{\partial \rho_{j-1}^n} ( \tilde{\rho}_{j-1}^n, \tilde{\rho}_{j}^n, \rho_{j+1}^n, \cdots, \rho_{j+m}^n)\geq 0,
        \quad \frac{\partial \mathcal{H}}{\partial \rho_{j}^n} ( \tilde{\rho}_{j-1}^n, \tilde{\rho}_{j}^n, \rho_{j+1}^n, \cdots, \rho_{j+m}^n)\geq 0.
    \]


    Since the derivative of $\mathcal{H}$ with respect to $\rho_{j-1}^n$ is given by
    \[
        \frac{\partial \mathcal{H}}{\partial \rho_{j-1}^n}(\rho_{j-1}^n, \rho_{j}^n, \rho_{j+1}^n, \cdots, \rho_{j+m}^n)
        =\lambda(1-q_j^n),
    \]
    where
    \[
        q_j^n = w_0 \rho_{j}^n + w_1 \rho_{j+1}^n+ \cdots + w_{m-1} \rho_{j+m-1}^n
        = \sum_{k=0}^{m-1} w_k \rho_{j+k}^n,
    \]
    we get
    \begin{align*}
        \frac{\partial \mathcal{H}}{\partial \rho_{j-1}^n} ( \tilde{\rho}_{j-1}^n, \tilde{\rho}_{j}^n, \rho_{j+1}^n, \cdots, \rho_{j+m}^n) = \lambda(1- \tilde{q}_j^n),
    \end{align*}
    where
    \[
        \tilde{q}_j^n
        = w_0 \tilde{\rho_{j}^n} + w_1 \rho_{j+1}^n+ \cdots + w_{m-1} \rho_{j+m-1}^n
        = w_0 \tilde{\rho_{j}^n}+ \sum_{k=1}^{m-1} w_k \rho_{j+k}^n.
    \]
    Moreover, we have $0 \leq \tilde{q}_j^n \leq 1$, because
    \[0 \leq \rho_{\min} \leq \tilde{\rho}_{j}^n \leq \rho_{\max} \leq 1 \qquad\text{and } \qquad 0 \leq \rho_{\min} \leq \rho_{j+k}^n \leq \rho_{\max} \leq 1\] imply
    \begin{align*}
        \tilde{q}_j^n & = w_0 \tilde{\rho_{j}^n}+ \sum_{k=1}^{m-1} w_k \rho_{j+k}^n \\
                      & \leq w_0 \rho_{\max} + \sum_{k=1}^{m-1} w_k \rho_{\max}     \\
                      & = \sum_{k=0}^{m-1} w_k \rho_{\max}                          \\
                      & = \rho_{\max} \sum_{k=0}^{m-1} w_k                          \\
                      & = \rho_{\max} \leq 1,
    \end{align*}
    where we have used the normalization condition \[\sum_{k=0}^{m-1} w_k = 1.\]
    By similar argument $\tilde{q}_j^n \geq \rho_{\min} \geq 0$. Therefore, we get
    \[\frac{\partial \mathcal{H}}{\partial \rho_{j-1}^n} ( \tilde{\rho}_{j-1}^n, \tilde{\rho}_{j}^n, \rho_{j+1}^n, \cdots, \rho_{j+m}^n) = \lambda(1- \tilde{q}_j^n) \geq 0.\]

    As the derivative of $\mathcal{H}$ with respect to $\rho_{j}^n$ is given by
    \[
        \frac{\partial \mathcal{H}}{\partial \rho_{j}^n}(\rho_{j-1}^n, \rho_{j}^n, \rho_{j+1}^n, \cdots, \rho_{j+m}^n)
        = 1-\lambda\left[(1-q_{j+1}^n)+w_0\rho_{j-1}^n\right],
    \]
    where
    \[
        q_{j+1}^n = w_0 \rho_{j+1}^n + w_1 \rho_{j+2}^n+ \cdots + w_{m-1} \rho_{j+m}^n = \sum_{k=0}^{m-1} w_k \rho_{j+k+1}^n,
    \]
    we get
    \[
        \frac{\partial \mathcal{H}}{\partial \rho_{j}^n} ( \tilde{\rho}_{j-1}^n, \tilde{\rho}_{j}^n, \rho_{j+1}^n, \cdots, \rho_{j+m}^n) =
        1-\lambda\left[(1-q_{j+1}^n)+w_0 \tilde{\rho}_{j-1}^n\right].\]
    We are looking to obtain
    \[
        \frac{\partial \mathcal{H}}{\partial \rho_{j}^n} ( \tilde{\rho}_{j-1}^n, \tilde{\rho}_{j}^n, \rho_{j+1}^n, \cdots, \rho_{j+m}^n) =
        1-\lambda\left[(1-q_{j+1}^n)+w_0\tilde{\rho}_{j-1}^n \right] \geq 0.\]
    In order to get this, we must have that
    \[
        0 \leq \lambda\left[(1-q_{j+1}^n)+w_0\tilde{\rho}_{j-1}^n\right]\leq 1
    \]
    Since we have $\lambda\left[(1-q_{j+1}^n)+w_0\rho_{j-1}^n\right] \geq 0$, it remains to show that
    \[\lambda\left[(1-q_{j+1}^n)+w_0\tilde{\rho}_{j-1}^n \right] \leq 1.\]
    We can estimate this as follows
    \begin{align*}
        \lambda\left[(1-q_{j+1}^n)+w_0\tilde{\rho}_{j-1}^n\right] & \leq \lambda\left[(1-q_{j+1}^n)+\tilde{\rho}_{j-1}^n\right]               \qquad \text{ since $0\leq w_0 \leq 1$ and $\tilde{\rho}_{j-1}^n \geq 0$ } \\
                                                                  & = \lambda \Bigl(\theta^{(1)}(q_{j}^{n}, q_{j+1}^{n}) - \theta^{(4)}(\tilde{\rho}_{j-1}^n,\tilde{\rho}_j^n)\Bigr)                                     \\
                                                                  & \leq \lambda \Bigl(\|\theta^{(1)}\|_{L^{\infty}}+ \|\theta^{(4)}\|_{L^{\infty}}\Bigr)                                                                \\
                                                                  & < 1, \qquad \text{ by Assumption 4.2.4 }.
    \end{align*}
    So far, we have shown that
    \[
        \frac{\partial \mathcal{H}}{\partial \rho_{j-1}^n} ( \tilde{\rho}_{j-1}^n, \tilde{\rho}_{j}^n, \rho_{j+1}^n, \cdots, \rho_{j+m}^n)(\rho_{j-1}^n- \rho_{\min}) \geq 0
    \]
    and
    \[
        \frac{\partial \mathcal{H}}{\partial \rho_{j}^n} ( \tilde{\rho}_{j-1}^n, \tilde{\rho}_{j}^n, \rho_{j+1}^n, \cdots, \rho_{j+m}^n)(\rho_{j}^n- \rho_{\min}) \geq 0.
    \]

    Therefore, $\Delta \mathcal{H}_1 \geq 0$.

    Now it remains to show $\Delta \mathcal{H}_2 \geq 0$. By the mean value theorem,
    there exist intermediate values
    \[ 0 \leq \rho_{\min} \leq \tilde{\rho}_{j+k}^n \leq \rho_{\max} \leq 1 \qquad \text{for } k= 1,2,\cdots, m,\]
    such that
    \[
        \Delta \mathcal{H}_2
        =
        \sum_{k=1}^{m}
        \frac{\partial \mathcal{H}}{\partial \rho_{j+k}^n}
        \bigl(
        \rho_{\min},\rho_{\min},
        \tilde\rho_{j+1}^n,\ldots,\tilde\rho_{j+m}^n
        \bigr)
        (\rho_{j+k}^n-\rho_{\min}). \]
    In order to have $\Delta \mathcal{H}_2 \geq 0$, we must have that
    \[
        \sum_{k=1}^{m}
        \frac{\partial \mathcal{H}}{\partial \rho_{j+k}^n}
        \bigl(
        \rho_{\min},\rho_{\min},
        \tilde\rho_{j+1}^n,\ldots,\tilde\rho_{j+m}^n
        \bigr)
        (\rho_{j+k}^n-\rho_{\min}) \geq 0.
    \]
    That means we need to show that
    \[
        \frac{\partial \mathcal{H}}{\partial \rho_{j+k}^n}
        \bigl(
        \rho_{\min},\rho_{\min},
        \tilde\rho_{j+1}^n,\ldots,\tilde\rho_{j+m}^n
        \bigr)
        (\rho_{j+k}^n-\rho_{\min}) \geq 0 \qquad \text{ for } k= 1,2, \cdots, m.
    \]
    Since
    \[
        (\rho_{j+k}^n-\rho_{\min}) \geq 0 \qquad \text{ for } k= 1,2, \cdots, m,
    \]
    it suffices to show that
    \[
        \frac{\partial \mathcal{H}}{\partial \rho_{j+k}^n}
        \bigl(
        \rho_{\min},\rho_{\min},
        \tilde\rho_{j+1}^n,\ldots,\tilde\rho_{j+m}^n
        \bigr)
        \geq 0 \qquad \text{ for } k= 1,2, \cdots, m.
    \]

    For $k=m$, the derivative of $\mathcal{H}$ with respect to $\rho_{j+m}^n$ is given by
    \[
        \frac{\partial \mathcal{H}}{\partial \rho_{j+m}^n}(\rho_{j-1}^n, \rho_{j}^n, \rho_{j+1}^n, \cdots, \rho_{j+m}^n)
        =
        \lambda w_{k-1}\rho_j^n.
    \]
    This leads us to
    \[
        \frac{\partial \mathcal{H}}{\partial \rho_{j+m}^n}\bigl(
        \rho_{\min},\rho_{\min},
        \tilde\rho_{j+1}^n,\ldots,\tilde\rho_{j+m}^n
        \bigr)
        =
        \lambda w_{k-1}\rho_{\min} \geq 0.
    \]
    Since the derivative of $\mathcal{H}$ with respect to $\rho_{j+k}^n$ is given by

    \[
        \frac{\partial \mathcal{H}}{\partial \rho_{j+k}^n}(\rho_{j-1}^n, \rho_{j}^n, \rho_{j+1}^n, \cdots, \rho_{j+m}^n)
        =
        \lambda\left(w_{k-1}\rho_j^n-w_k\rho_{j-1}^n\right),
        \qquad \text{ for } k=1,\ldots,m,
    \]
    we get
    \begin{align*}
        \frac{\partial \mathcal{H}}{\partial \rho_{j+k}^n}
        \bigl(
        \rho_{\min},\rho_{\min},
        \tilde{\rho}_{j+1}^n,\ldots,\tilde\rho_{j+m}^n
        \bigr) & = \lambda \left(w_{k-1}\rho_{\min} - w_{k}\rho_{\min}\right) \\
               & = \lambda \rho_{\min} (w_{k-1}- w_{k})
    \end{align*}
    By Assumption 4.2.3 we have $(w_{k-1}- w_{k}) \geq cm^{-2} > 0$. Therefore,
    \[
        \frac{\partial \mathcal{H}}{\partial \rho_{j+k}^n}
        \bigl(
        \rho_{\min},\rho_{\min},
        \tilde\rho_{j+1}^n,\ldots,\tilde\rho_{j+m}^n
        \bigr)= \lambda \rho_{\min} \left(w_{k-1}- w_{k}\right) \geq 0.
    \]
    So, it follows from that we obtain
    \[
        \Delta \mathcal{H}_2
        =
        \sum_{k=1}^{m}
        \frac{\partial \mathcal{H}}{\partial \rho_{j+k}^n}
        \bigl(
        \rho_{\min},\rho_{\min},
        \tilde\rho_{j+1}^n,\ldots,\tilde\rho_{j+m}^n
        \bigr)
        (\rho_{j+k}^n-\rho_{\min}) \geq 0\]

    We have shown that both $ \Delta \mathcal{H}_1 \geq 0 $ and $ \Delta \mathcal{H}_2 \geq 0$. Which implies
    \[
        \mathcal{H}(\rho_{j-1}^n,\rho_j^n,\rho_{j+1}^n,\ldots,\rho_{j+m}^n)-\rho_{\min}
        =
        \Delta \mathcal{H}_1+\Delta \mathcal{H}_2 \geq 0 ,\]
    and therefore
    \[
        \mathcal{H}(\rho_{j-1}^n,\rho_j^n,\rho_{j+1}^n,\ldots,\rho_{j+m}^n) \geq \rho_{\min}.
    \]
    The upper bound can be shown by similar argument.
\end{proof}
Now we proof the Maximum principle by induction.
\begin{proof}[Proof of Theorem 4.2.7]
    For $n=0$, it is obvious that $\rho_{\min} \leq \rho_j^0 \leq \rho_{\max}$ for all $ j \in \mathbb{Z}$. Assume now that the result holds for any $n \in \mathbb{N}$, that is,
    \[
        \rho_{\min} \leq \rho_j^n \leq \rho_{\max}  \qquad \text{for all } j \in \mathbb{Z}.
    \] We show that it also holds for $n+1$. By Lemma 4.2.8, we have
    \[
        \rho_{\min} \leq \mathcal{H}\bigl(\rho_{j-1}^n,\rho_j^n, \rho_{j+1}^n, \ldots,\rho_{j+m}^n \bigr) \leq
        \rho_{\max}.
    \]
    By the definition of $\rho_j^{n+1}$, this gives $\rho_{\min} \leq \rho_j^{n+1} \leq \rho_{\max}$. Thus, by induction,
    \[
        \rho_{\min} \leq \rho_j^n \leq \rho_{\max} \qquad \text{for all }  j \in \mathbb{Z} \quad \text{and } n \in \mathbb{N}.
    \]
\end{proof}
\subsection{One-sided Lipschitz Estimate}
Next, we are going to prove the one-sided Lipschitz condition. The one-sided Lipschitz condition will play a crucial role in showing the monotonicity of the update function $\mathcal{H}$ given in \eqref{eq:mono_H}.
\begin{lemma}
    Assume all conditions in Theorem 4.2.7 are given. Let $\{\rho_j^n\}_{j \in \mathbb{Z}}^{ n \in \mathbb{N}}$ be the numerical solution computed by \eqref{eq:mono_H}. Then, the following numerical differences
    \begin{align*}
        r_j^n = \rho_{j+1}^n - \rho_j^n, \qquad j \in \mathbb{Z}, n \in \mathbb{N},
    \end{align*}
    satisfy the one-sided Lipschitz condition
    \begin{equation}\label{eq:one_sided_condition}
        r_j^n \geq -L\Delta x,  \qquad j \in \mathbb{Z}, n \in \mathbb{N}.
    \end{equation}
\end{lemma}
Before providing the proof of the above Lemma, let us first study the structure of
\[
    r_j^{n+1} =  \rho_{j+1}^{n+1} - \rho_j^{n+1}.
\]
Since we have
\[
    \rho_{j+1}^{n+1} = \rho_{j+1}^{n} + \lambda\Bigl[ g(\rho_{j}^{n}, \rho_{j+1}^{n}, q_{j}^{n},q_{j+1}^{n})-g(\rho_{j+1}^{n}, \rho_{j+2}^{n}, q_{j+1}^{n},q_{j+2}^{n})\Bigr]
\]
and
\[
    \rho_{j}^{n+1} = \rho_{j}^{n} + \lambda\Bigl[ g(\rho_{j-1}^{n}, \rho_{j}^{n}, q_{j-1}^{n},q_{j}^{n})-g(\rho_{j}^{n}, \rho_{j+1}^{n}, q_{j}^{n},q_{j+1}^{n})\Bigr],
\]
taking the the difference gives us
\begin{equation*}
    \begin{split}
        r_j^{n+1} & = \rho_{j+1}^{n+1} - \rho_j^{n+1}                                                                                                                                                                                            \\
                  & = (\rho_{j+1}^n - \rho_j^n) + \lambda\Bigl[ 2g(\rho_{j}^{n}, \rho_{j+1}^{n}, q_{j}^{n},q_{j+1}^{n})-g(\rho_{j-1}^{n}, \rho_{j}^{n}, q_{j-1}^{n},q_{j}^{n})- g(\rho_{j+1}^{n}, \rho_{j+2}^{n}, q_{j+1}^{n},q_{j+2}^{n})\Bigr] \\
                  & =  r_j^n +  \lambda\Bigl[ 2g(\rho_{j}^{n}, \rho_{j+1}^{n}, q_{j}^{n},q_{j+1}^{n})-g(\rho_{j-1}^{n}, \rho_{j}^{n}, q_{j-1}^{n},q_{j}^{n})- g(\rho_{j+1}^{n}, \rho_{j+2}^{n}, q_{j+1}^{n},q_{j+2}^{n})\Bigr].
    \end{split}
\end{equation*}
As $g$ is quadratic, Huang and Du, see \cite[p.~12]{huangAsymptoticCompatibilityClass2024}, used Taylor's expansions to write $r_j^{n+1}$ as following:
\begin{equation}\label{eq:differences}
    r_j^{n+1}
    =
    \sum_{k=-1}^m c_{j,k}^n r_{j+k}^n
    -2\lambda(\gamma_{13}+\gamma_{14})(r_j^n)^2
    -2\lambda(\gamma_{23}+\gamma_{24})(r_{j+1}^n)^2,
\end{equation}
such that
\[\sum_{k=-1}^m c_{j,k}^n = 1.\]

Since we have
\[ \gamma_{13}= \gamma_{23}= \gamma_{24}=0 \qquad \text{and } \gamma_{14}\]
for the Godunov numerical flux $g$, \eqref{eq:differences} simplifies to
\begin{equation}\label{eq:differences_God}
    r_j^{n+1}
    =
    \sum_{k=-1}^m c_{j,k}^n r_{j+k}^n
    +2\lambda(r_j^n)^2,
\end{equation}
where,
\begin{align*}
    c_{j,-1}^n & = \lambda(1-q_j^n)                                           \\
    c_{j,0}^n  & = 1 - \lambda (1- q_j^n) + \lambda p_{j,0}^n -2\lambda r_j^n \\
    c_{j,k}^n  & = \lambda p_{j,k}^n, \qquad \text{for } k= 1,2,\cdots,m,
\end{align*}
and
\[ p_{j,k}^n =  w_{k-1}\rho_{j+1}^n - w_k\rho_j^n + w_k r_j^n \qquad \text{with convention } w_{-1}= w_m = 0 \]

Now we prove the Lemma by induction
\begin{proof}[Proof of Lemma 4.2.9]
    The one-sided Lipschitz condition on $\rho_0$ in Assumption 4.2.2 gives us $r_j^0 \geq - L\Delta x$ for all $j\in \mathbb{Z}$. Assume now that \eqref{eq:one_sided_condition} holds for any $n \in \mathbb{N}$, that is,
    \[
        r_j^n \geq - L\Delta x,  \qquad j \in \mathbb{Z}, n \in \mathbb{N}.
    \]
    We show also \eqref{eq:one_sided_condition} holds for $n+1$.
    It suffices to show that $r_j^{n+1}$ given in \eqref{eq:differences_God} is a convex combination of $r_{j-1}^n, r_{j}^n,\cdots,r_{j+m}^n$ plus the nonnegative term $2\lambda(r_j^n)^2$, since this implies
    \[ \inf\limits_{j \in \mathbb{Z}} r_j^{n+1} \geq \inf\limits_{j \in \mathbb{Z}} r_j^{n} \geq -L\Delta x.\]
    As the $\lambda(r_j^n)^2$ is nonnegative and $\sum_{k=-1}^m c_{j,k}^n = 1$, it remains to show that $c_{j,k}^n \geq 0$ for all $k= -1, 0,1, \cdots, m.$

    Clearly, $c_{j,-1}^n \geq 0$ because we have shown that in \eqref{eq:First_argument}
    \begin{align*}
        c_{j,-1}^n & =\frac{\partial \mathcal{H}}{\partial \rho_{j-1}^n}              \\
                   & = \lambda(1-q_j^n) \geq 0 \qquad \text{as } 0 \leq q_j^n \leq 1.
    \end{align*}
    For $k= 0$, we have
    \begin{align*}
        c_{j,0}^n & = 1 - \lambda (1- q_j^n) + \lambda p_{j,0}^n -2\lambda r_j^n                            \\
                  & = 1 - \lambda (1- q_j^n) + \lambda \Bigl(-w_0\rho_j^n + w_0 r_j^n\Bigr) -2\lambda r_j^n \\
                  & = 1- \lambda \Bigl[(1- q_j^n)+ w_0\rho_j^n  + (2-w_0)r_j^n\Bigr]
    \end{align*}
    In order to have  $c_{j,0}^n \geq 0$, it holds to show that $\lambda \Bigl[(1- q_j^n)+ w_0\rho_j^n  + (2-w_0)r_j^n\Bigr] \leq 1$. But by \eqref{eq:lambda_God} we have $\lambda < \frac{1}{2}$. Therefore, it is sufficient to show that
    \[
        (1- q_j^n)+ w_0\rho_j^n  + (2-w_0)r_j^n \leq 2.
    \]
    Mark that all the terms $(1- q_j^n)$, $w_0\rho_j^n$ and $(2-w_0)$ are nonnegative provided that $ 0\leq q_j^n \leq 1$,  $0 \leq  w_0 \leq 1$ and  $0\leq \rho_j^n \leq 1$.

    Now consider that  $r_j^n \leq 0$, then we $(2-w_0)r_j^n \leq 0$ and
    \begin{align*}
        (1- q_j^n)+ w_0\rho_j^n  + (2-w_0)r_j^n & \leq (1- q_j^n)+ w_0\rho_j^n \\
                                                & \leq (1- q_j^n)+ \rho_j^n    \\
                                                & \leq 2
    \end{align*}
    since  $ 0\leq q_j^n \leq 1$, $0 \leq  w_0 \leq 1$, and  $0\leq  \rho_j^n \leq 1$.
\end{proof}

%\[ \tilde{q}_j^n = \sum_{k=0}^{m-1} \tilde{\rho}_{j+k}^n\]
%\section{Assumptions}
%\section{AC Theorem}
%\section{Maximum Principle}

%\section{TVD in Space and Temporal TV-estimate}
